\section{次数一から六の全根共鳴と混合simple-root locus}
\label{sec:maximal-all-degrees}
前節のcubic（$E_6$）の構成は，$r=3,\ldots,8$点の
$\mathbb P^2$のblow-upに一様に拡張する．
Rains\cite[\S4, (4.15)]{RainsFiltered}の根系は
\[
 \begin{array}{c|rrrrrr}
 r&3&4&5&6&7&8\\\hline
 d=9-r&6&5&4&3&2&1\\
 \text{根系}&A_2+A_1&A_4&D_5&E_6&E_7&E_8\\
 \#\text{根}&8&20&40&72&126&240
 \end{array}
\]
である．

\begin{theorem}[全次数における最大根共鳴の稠密性]
\label{thm:all-degree-maximal}
$E\subset\mathbb P^2$を変曲点$o$を原点とする滑らかな
楕円三次とし，$q\in E$を非torsionとする．
$3\le r\le8$として，互いに異なる正整数
$N_1,\ldots,N_r$を選び，$p_i=N_iq$と置く．
すると
\[
 S_r=\operatorname{Bl}_{p_1,\ldots,p_r}\mathbb P^2
\]
は$-K_{S_r}$がampleな滑らかなgenuine del Pezzo曲面
（次数$d=9-r$）である．また，
\emph{根系の全ての根}に対して
\[
 \rho(\alpha)=q^{-m(\alpha)},\qquad m(\alpha)\ne0
\]
が成立する．このような三以上八以下の点配置は，
固定した$E,q$に対する$E^r$にZariski稠密である．
$N_1<N_3<N_2$という条件を加えても同様に稠密である．

根系の任意の根が共鳴するという条件は，
六点の場合と同じ単純根
\[
 \alpha_0=h-e_1-e_2-e_3,\quad
 \alpha_i=e_i-e_{i+1}\quad(1\le i\le r-1)
\]
だけについて共鳴することと同値であり，
$r$個の整数指数$m_0,\ldots,m_{r-1}$による
\begin{equation}\label{eq:all-degree-chart}
 3p_1=(m_0-2m_1-m_2)q,\qquad
 p_i=p_1+\bigl(m_1+\cdots+m_{i-1}\bigr)q
 \quad(2\le i\le r)
\end{equation}
で完全に記述される．$p_1$の選択は$E[3]$の九通り．
このmaximal resonance locusで$S_r$がgenuineとなることは
任意の根の指数が零でないことと同値である．
\end{theorem}
\begin{proof}
根は$a h-\sum_{i=1}^r b_ie_i$の形で
\[
 \sum b_i=3a,\qquad \sum b_i^2=a^2+2
\]
を満たす．Cauchy--Schwarz不等式から
$(9-r)a^2\le2r$なので$|a|\le4$．
$a\ge0$として分類すれば，
$a=0$では$e_i-e_j$，
$a=1$では$h-e_i-e_j-e_k$，
$a=2$では$2h-e_{i_1}-\cdots-e_{i_6}$が出る．
$r=8$の$a=3$に限って
$3h-2e_i-\sum_{j\ne i}e_j$も現れる．
これらの反数も根であり，上の表の根数を尽くす．

ここで整数$b$について$b(b-1)\ge0$であり，
\[
 \sum_i b_i(b_i-1)=(a^2+2)-3a
\]
である．$a=1,2$では右辺が零なので全$b_i$は$0$または$1$，
$a=3$では右辺が$2$なので，$r=8$で一つの$b_i=2$，
残る七つは$1$である．$r=8,a=4$なら右辺は$6$で，
一つの$3$（または$-2$）と残る$0,1$，
あるいは三つの$2$（または$-1$）と残る$0,1$
しか可能でないが，いずれも$\sum b_i=12$に到達しない．
従って$a=4$の根は存在しない．

$p_i=N_iq$なら，$\rho(a h-\sum b_ie_i)=q^{-\sum b_iN_i}$．
$a=0$の根では指数は相異なる二整数の差，
$a>0$の上の根では指数は正整数の正係数付き和となる．
従って全ての根で$m(\alpha)\ne0$である．
$S_r$上では滑らかな反標準曲線$E'$が存在し，
$H=-K_{S_r}=[E']$は$H^2=9-r>0$でnefである：
$H$と任意の相異なる既約曲線の交点は非負だからである．
もし$H$がampleでなければ，nefかつbigの曲面版
Nakai--Moishezon判定により$H\cdot C=0$なる既約曲線$C$
がある．Hodge指数定理により$C^2<0$，
adjunctionから$C^2=2p_a(C)-2$なので
$C^2=-2$，$p_a(C)=0$．
$C$と$E'$は交わらないから，$C|_{E'}$は自明であり，
その根類の制限は$\mathcal O_E$となる．
これは$m(C)\ne0$に矛盾する．よって$H$はample．

根の整数表示\eqref{eq:all-degree-chart}は
前節の定理\ref{thm:maximal-resonance}と全く同じ
単純根による基底表示から従う．
逆に$m(\alpha)=0$なる根があると
その制限は自明となる．点配置がsmooth genuineなら
このような根は存在しない（命題\ref{prop:zero-index}）．
したがってgenuine条件はregular index条件に同値．
最後のZariski稠密性と大小制約の稠密性も，
前節の系\ref{cor:maximal-dense}と同じ
正整数倍の直積および$E^3$のユニモジュラ自己同型の議論による．
\end{proof}

\begin{corollary}[各次数の実在混合根と反標準AZ障害]
\label{cor:all-degree-mixed}
各$d=1,\ldots,6$のRains型smooth elliptic
noncommutative del Pezzoの当該blow-up chartで，
$N_1<N_3<N_2$を満たす上記のZariski稠密な
maximal resonance点では，単純根
$\alpha=e_1-e_2$，$\beta=e_2-e_3$の指数は
$m_\alpha=N_2-N_1>0$，
$m_\beta=N_3-N_2<0$である．
Rainsのsimple-root effectivity・既約sheaf存在結果が
適用されると，$\alpha,\beta$に対応する
境界非交叉の既約純一次元sheafが同じRains heartに
同時に存在し，$\dim\Ext^1(M_\alpha,M_\beta)=1$となる．
特に$M_\beta$の存在により$(e,\tau)$は
Rains heartでArtin--Zhang ampleでない．

この結論はそのRains elliptic chartの命題であり，
全非可換del Pezzoモジュライの全既約成分の
Zariski稠密な障害を主張したものではない．
\end{corollary}
\begin{proof}
$\alpha\cdot\beta=1$，各指数の符号は直接計算で得られる．
Rains\cite[\S4]{RainsFiltered}のsimple-root
effectivity判定と\cite[\S\S9,11]{RainsBirational}を使う．
残りは命題\ref{prop:root-intersection}と
定理\ref{thm:sign}から従う．
\end{proof}
